% Abhakseite „Das kann ich“ – Zone nur im Gesamt (\mitzone)
%% TODO Ich-kann-Sätze: Platzhalter wie in den Aufgabendateien
\begin{abhakseite}
\ifdefined\mitzone
\abhakgruppe{Kennst du schon}
\abhak{1}{Baumdiagramm und Pfadregeln mit Zurücklegen}
\abhak{2}{Ziehen mit und ohne Zurücklegen unterscheiden (Urnenmodell)}
\abhak{3}{Binomialkoeffizient (n über k) als Anzahl der Auswahlen von k aus n, mit der Rechnertaste nCr und für kleine Zahlen im Kopf; Fakultät}
\abhak{4}{Rechnerfunktionen für die Binomialverteilung (Einzelwahrscheinlichkeit und kumulierte Wahrscheinlichkeit, je nach Gerät binompdf und binomcdf) bedienen oder die Tabelle der summierten Binomialverteilung lesen (Spalte p, Zeile k; bei p über ein Halb Treffer und Niete tauschen)}
\abhak{5}{Anteile, Prozente und Anzahlen ineinander umrechnen und eine Ungleichung nach der gesuchten Größe umformen (Ungleichheitszeichen dreht beim Teilen durch eine negative Zahl), Ergebnis nach dem Sinn auf- oder abrunden}
\abhak{6}{Säulendiagramm lesen (Höhe am Gitter ablesen, Säulen addieren) und zeichnen; Summenzeichen lesen (untere und obere Grenze, Laufvariable)}
\abhak{7}{Potenzen mit einer Basis zwischen null und eins}
\abhak{8}{Erwartungswert μ = n · p berechnen und als mittlere Trefferzahl deuten}
\abhak{9}{Rechnerfunktionen für die Binomialverteilung (Einzelwahrscheinlichkeit und kumulierte Wahrscheinlichkeit, je nach Gerät binompdf und binomcdf) bedienen oder die Tabelle der summierten Binomialverteilung lesen (Spalte p, Zeile k; bei p über ein Halb Treffer und Niete tauschen) – Fehler finden}
\abhak{10}{Rechnerfunktionen für die Binomialverteilung (Einzelwahrscheinlichkeit und kumulierte Wahrscheinlichkeit, je nach Gerät binompdf und binomcdf) bedienen oder die Tabelle der summierten Binomialverteilung lesen (Spalte p, Zeile k; bei p über ein Halb Treffer und Niete tauschen) – selbst rechnen}
\fi
\abhakgruppe{Einheit 1 · Bernoulli-Experiment und Bernoulli-Kette}
\abhak{11}{Modell erkennen}
\abhak{12}{Modell erkennen – Fehler finden, Begründen, Anwendung, Darstellungswechsel}
\abhakgruppe{Einheit 2 · Bernoulli-Formel}
\abhak{13}{Bernoulli-Formel}
\abhak{14}{Bernoulli-Formel – Fehler finden, Begründen, Anwendung, Darstellungswechsel}
\abhakgruppe{Einheit 3 · Kumulierte Wahrscheinlichkeiten}
\abhak{15}{Kumulierte Wahrscheinlichkeiten}
\abhak{16}{Bedingte Restwahrscheinlichkeit nach bekannten Ergebnissen über die Binomialverteilung berechnen}
\abhak{17}{Wahrscheinlichkeit für zwei unabhängige Spieler als Produkt binomialer Wahrscheinlichkeiten berechnen}
\abhak{18}{Kumulierte Wahrscheinlichkeiten – Fehler finden, Begründen, Anwendung, Darstellungswechsel}
\abhakgruppe{Einheit 4 · Umkehraufgaben}
\abhak{19}{Umkehraufgaben}
\abhak{20}{Trefferzahlen mit einer Einzelwahrscheinlichkeit über einer Schranke mit dem Rechner ermitteln}
\abhak{21}{Umkehraufgaben – Fehler finden, Begründen, Anwendung}
\abhakgruppe{Einheit 5 · Verteilung im Diagramm}
\abhak{22}{Verteilung im Diagramm}
\abhak{23}{Obere Grenze einer im Diagramm markierten kumulierten Wahrscheinlichkeit über den Erwartungswert ermitteln}
\abhak{24}{Verteilung im Diagramm – Fehler finden, Begründen, Anwendung, Darstellungswechsel}
\end{abhakseite}
